% ============================================================
% Weekly Meeting Report — Week 24
% Topic: Supervisor Feedback on Draft; Per-Layer Drift Analysis
% ============================================================

\section*{Supervisor Feedback \& Response Analysis}
\addcontentsline{toc}{section}{Supervisor Feedback \& Response Analysis}

\begin{tabular}{ll}
    \textbf{Date}       & 16~June~2026 \\
    \textbf{Student}    & Khac Duc Giang Nguyen \\
    \textbf{Supervisor} & Dr.\ Seyed Sahand Mohammadi Ziabari \\
    \textbf{Phase}      & Post-draft review; analysis jobs; response to supervisor \\
\end{tabular}

\vspace{1em}

% ------------------------------------------------------------
\subsection*{1. Overview}

This report covers the period 5--16~June~2026 (Week~24). The supervisor
returned written feedback on the thesis draft, raising six points requiring
deeper investigation or stronger evidence. This report documents each point,
our analysis of what can be answered from existing data, what requires new
experiments, and what goes into the limitations section. Two additional
analysis jobs were submitted on Snellius (jobs 23937817 and the existing
eval output from 23821658 was re-examined). Total additional compute used:
$\approx$1~SBU.

% ------------------------------------------------------------
\subsection*{2. Supervisor Feedback --- Full Text}

The supervisor's email, received 16~June~2026, raised six points:

\begin{enumerate}
    \item \textbf{Gradient starvation hypothesis is indirect.}
    High cosine similarity alone does not prove starvation over localised
    changes in a few critical layers. Have gradients been measured by target
    stratum directly, or is the claim inferred from dataset composition?

    \item \textbf{Large-target collapse after one epoch.}
    Could this reflect calibration failure rather than forgetting? Recall,
    confidence distributions, and false positives should be analysed separately.

    \item \textbf{Normal stratum also collapses despite CST containing
    some normal targets.}
    Why is the remaining gradient insufficient to preserve those representations?

    \item \textbf{Stage~2 KD / detection ratio increasing.}
    The rising kd/det ratio suggests distillation increasingly dominates
    optimisation. Without a $\lambda$ ablation (0, 0.33, 1.0) or scheduled
    $\lambda_{\mathrm{kd}}$, it is unclear whether KD constrains adaptation
    too strongly.

    \item \textbf{Stage~2 tiny AP $\approx$ 0.}
    If the model never learned tiny targets before Stage~3, is Stage~3 tiny
    collapse measuring forgetting or failure to learn an entirely new capability?

    \item \textbf{Central claim depends on separating gradient starvation
    from broader distribution shift.}
    How can missing-scale gradients be distinguished from changes in
    background statistics, thermal characteristics, or annotation noise?
\end{enumerate}

% ------------------------------------------------------------
\subsection*{3. Compute Remaining \& New Jobs Submitted}

\begin{table}[h]
\centering
\small
\begin{tabular}{ll}
\hline
Budget at Meeting~5      & $\approx$6,860 SBU remaining \\
Job 23937817 (drift\_s2s3, rome CPU) & $\approx$1 SBU \\
\textbf{Remaining after this session} & \textbf{$\approx$6,859 SBU} \\
\hline
\end{tabular}
\end{table}

\noindent Two analysis tasks were identified:
\begin{enumerate}
    \item \texttt{run\_drift\_s2s3.sh} --- per-layer-group parameter drift
    S2$\to$S3 using \texttt{parameter\_drift.py}. Submitted to the \texttt{rome}
    CPU partition. Completed in 25~seconds (job~23937817).
    \item PR curves and recall distributions for T1~val after Stage~3 ---
    these already existed from job~23821658 at
    \texttt{/projects/prjs2041/analysis/stage3\_naive/eval\_t1/}.
    No new job required.
\end{enumerate}

% ------------------------------------------------------------
\subsection*{4. Point-by-Point Response}

% ----
\subsubsection*{4.1 Gradient Starvation Is Indirect (Point~1)}

\textbf{Status: Partially addressed --- strengthened by new data, structural
argument added.}

\vspace{0.5em}
\noindent\textbf{New evidence from job~23937817.} \texttt{parameter\_drift.py}
was run comparing Stage~2 seed-42 best.pt to Stage~3 naive best.pt. 335 parameters
analysed. Summary for learnable parameters only (BN running stats excluded):

\begin{table}[h]
\centering
\caption{Per-layer-group parameter drift S2$\to$S3 (learnable params only).}
\small
\begin{tabular}{lrrrr}
\hline
\textbf{Group} & \textbf{N} & \textbf{Cos mean} & \textbf{Cos min} & \textbf{L2 mean} \\
\hline
backbone & 36 & 0.9951 & 0.9626 & 0.069 \\
head     &  3 & 0.9894 & 0.9697 & 0.100 \\
other    & 91 & 0.9911 & 0.9596 & 0.102 \\
neck     & 75 & 0.9772 & 0.8683 & 0.171 \\
\hline
\end{tabular}
\end{table}

\begin{figure}[h]
\centering
\begin{minipage}{0.48\linewidth}
    \includegraphics[width=\linewidth]{media/meetings/fig_drift_groups.png}
    \caption{Per-layer-group cosine similarity S2$\to$S3. Neck adapts most; head is effectively frozen.}
\end{minipage}\hfill
\begin{minipage}{0.48\linewidth}
    \includegraphics[width=\linewidth, height=0.35\textheight, keepaspectratio]{media/meetings/fig_drift_cos.png}
    \caption{Per-layer cosine similarity dot plot. No layer shows localised large change.}
\end{minipage}
\end{figure}

\noindent\textbf{Interpretation.}
\begin{itemize}
    \item No layer group shows substantial change: all learnable-parameter
    cosine similarities are $>$0.96. The supervisor's concern about
    ``localised changes in a few critical layers'' is answered empirically:
    no such layer exists. The lowest cosine similarity in the head group
    is 0.9697 (\texttt{model.24.conv.weight}).
    \item \textbf{Backbone is the most stable} (cos~0.9951) --- feature
    extraction is essentially unchanged. The model's visual representation
    of large targets still exists in the weights.
    \item \textbf{Neck changed the most} (cos~0.9772, min~0.8683) --- the
    feature pyramid network adapted to encode CST's tiny/small distribution.
    This shifts the feature space the detection head receives, creating a
    mismatch for large-target representations that the head never adapted to.
    \item \textbf{Head barely moved} (cos~0.9894) --- the detection prediction
    layers are effectively frozen. This is precisely what gradient starvation
    predicts: 0\% large targets in CST $\to$ zero positive gradient flows into
    the large-anchor detection outputs $\to$ those weights never update.
\end{itemize}

\noindent\textbf{Structural argument (not just empirical).}
In YOLOMG with P3/P4/P5 multi-scale outputs, large targets are detected
by the P5 branch (stride~32, coarsest grid). CST Anti-UAV has 0\% large
targets (longest bounding box side $\geq$64\,px). During Stage~3 training,
the detection loss contribution from large-target boxes at the P5 output is
identically zero --- not small, but mathematically zero. This is not a
hypothesis about gradient flow; it is a direct consequence of the detection
loss function~\cite{pezeshki2021gradient}. The empirical finding (head barely
moved, cos~0.9894) is consistent with and provides supporting evidence for
this structural argument. Masip et al.~\cite{masip2026face}
characterise catastrophic forgetting mechanistically as feature
\emph{rotation} or capacity reduction rather than bulk weight overwriting ---
the high per-layer cosine similarities here are precisely what that framework
predicts when gradient absence (rather than gradient interference) is the
operative mechanism.

\vspace{0.5em}
\noindent\textbf{What is still missing.} Actual gradient magnitudes per layer per
target stratum were not logged during Stage~3 training. This would require
re-running Stage~3 with per-batch gradient hooks --- not feasible within the
compute budget. This limitation is acknowledged in the thesis.

\vspace{0.5em}
\noindent\textbf{Thesis update.} The per-layer-group drift table and
\texttt{fig\_drift\_groups.png} should be added to Chapter~4 (Results),
alongside the S1$\to$S2 comparison. The argument should be reframed:
per-layer analysis rules out localised overwriting; structural analysis
explains why large-target gradient starvation is mechanically inevitable;
direct gradient measurement is left as future work.

% ----
\subsubsection*{4.2 Large-Target Collapse: Calibration or Forgetting? (Point~2)}

\textbf{Status: Fully answered from existing data --- no new experiments needed.}

\vspace{0.5em}
\noindent The PR curves and confidence distributions from job~23821658 already
exist locally at \texttt{analysis/stage\_3/stage3\_naive/eval\_t1/}. Key numbers
from \texttt{metrics\_summary.txt}:

\begin{table}[h]
\centering
\caption{T1 val evaluation after Stage~3 naive (job~23821658).}
\small
\begin{tabular}{lr}
\hline
\textbf{Metric} & \textbf{Value} \\
\hline
Overall mAP@0.5          & 0.0761 \\
Overall Precision        & 0.1342 \\
Overall Recall           & 0.2521 \\
Overall F1               & 0.1752 \\
Tiny mAP@0.5             & 0.0004 \\
Small mAP@0.5            & 0.0670 \\
Normal mAP@0.5           & 0.0878 \\
Large mAP@0.5            & \textbf{0.0000} \\
\hline
\end{tabular}
\end{table}

\begin{figure}[h]
\centering
\begin{minipage}{0.48\linewidth}
    \includegraphics[width=\linewidth]{media/meetings/fig_pr_curve.png}
    \caption{PR curve on T1 val after Stage~3. Large-stratum curve is absent (AP\,=\,0).}
\end{minipage}\hfill
\begin{minipage}{0.48\linewidth}
    \includegraphics[width=\linewidth]{media/meetings/fig_conf_curve.png}
    \caption{P/R/F1 vs confidence threshold. No operating point recovers large-target recall.}
\end{minipage}
\end{figure}

\begin{figure}[h]
\centering
\includegraphics[width=0.6\linewidth]{media/meetings/fig_scale_metrics.png}
\caption{Per-stratum mAP@0.5 after Stage~3. Large stratum: 0.000; small stratum remains non-zero (0.067).}
\end{figure}

\noindent\textbf{Why large mAP\,=\,0.000 rules out calibration failure.}
Average Precision (AP) is computed by integrating the precision-recall
curve across \textit{all} confidence thresholds (using the 101-point
interpolation). AP\,=\,0.000 means that at \textit{every} confidence
threshold, either recall\,=\,0 (no large targets detected) or
precision\,=\,0 (every detection is a false positive). Calibration failure
means confidence is miscalibrated but detection still occurs --- it would
manifest as non-zero recall with shifted precision. AP\,=\,0 rules this
out definitively. The model detects no large targets under any operating
point. This is true forgetting, not miscalibration.

\vspace{0.5em}
\noindent The \texttt{fig\_conf\_curve.png} (P/R/F1 vs confidence threshold)
and \texttt{fig\_pr\_curve.png} are available and can be included in the
thesis as supporting evidence.

\vspace{0.5em}
\noindent\textbf{Thesis update.} Add one sentence: ``Large-stratum AP\,=\,0.000
is computed by integrating over all confidence thresholds; miscalibrated
confidence would produce non-zero recall at some threshold. The zero-AP
result confirms complete loss of large-target detection rather than
confidence drift.'' Include \texttt{fig\_conf\_curve.png} in the appendix.

% ----
\subsubsection*{4.3 Normal Stratum Also Collapses (Point~3)}

\textbf{Status: Addressable by argument; exact CST normal percentage
should be confirmed from dataset audit.}

\vspace{0.5em}
\noindent CST Anti-UAV scale distribution: 67.3\,\% tiny, 30.4\,\% small,
$\approx$2.3\,\% normal, 0\,\% large (from earlier dataset audit). The
2.3\,\% normal fraction means normal-target gradient signal exists but is
severely imbalanced relative to Stage~1 training (where normal comprised
$\approx$60\,\% of targets).

\vspace{0.5em}
\noindent\textbf{Why the remaining gradient is insufficient:}
\begin{enumerate}
    \item \textbf{Gradient imbalance.} During each Stage~3 batch, $\approx$97.7\,\%
    of gradient signal comes from tiny/small targets and $\approx$2.3\,\% from
    normal targets. The optimiser update is dominated by tiny/small loss terms.
    Normal-target representations are pushed toward a compromise that serves
    the dominant distribution~\cite{molahasani2023continual,he2024gradientreweighting}.
    \item \textbf{Feature distribution shift in the neck.} The neck adapts to
    CST's tiny/small dominant distribution (cos~0.9772, most changed of all
    groups in the per-layer analysis). Even if the head's normal-anchor weights
    were preserved, the neck now maps normal-sized CST inputs through a feature
    space optimised for tiny/small detection. This breaks the feature-to-detection
    correspondence for normal-sized Stage~1 targets, even though some normal
    gradient does flow.
    \item \textbf{Domain gap within the normal stratum.} CST normal targets
    (2.3\,\%) likely differ in thermal appearance, background clutter, and
    acquisition distance from Anti-UAV-RGBT normal targets ($\approx$60\,\%).
    The gradient from CST normal instances may update normal-target representations
    toward a CST-specific distribution rather than preserving the Stage~1
    distribution.
\end{enumerate}

\noindent\textbf{Thesis update.} Add the gradient imbalance framing:
``Although CST contains $\approx$2.3\,\% normal-sized targets, the
gradient signal from these instances ($\approx$2.3\,\% of batch
gradient) is insufficient to counteract two effects: the dominant
tiny/small gradient pushing the feature pyramid toward smaller-scale
representations, and the thermal domain gap between CST normal
instances and Anti-UAV-RGBT normal instances.''

% ----
\subsubsection*{4.4 $\lambda_{\mathrm{kd}}$ Ablation Not Performed (Point~4)}

\textbf{Status: Genuine gap --- goes in limitations.}

\vspace{0.5em}
\noindent Three Stage~2 seed runs all used $\lambda_{\mathrm{kd}}\,=\,1.0$.
The single-seed pilot (job~22688135, A100 partition, FM\,=\,$-$0.0153) used
a different architecture and DDP setup and is not a clean $\lambda\,=\,0$
baseline. No controlled $\lambda$ ablation (0, 0.33, 1.0) was run.

\vspace{0.5em}
\noindent\textbf{What the data does show:}
\begin{itemize}
    \item The kd/det ratio rises from 1.36$\times$ at epoch~0 to 2.44$\times$
    at epoch~16 (best checkpoint). KD increasingly dominates optimisation.
    \item FM\,=\,$-$0.033\,$\pm$\,0.004 is the empirical result under this
    regime. Whether a lower $\lambda$ would achieve similar T1 retention
    with better T2 adaptation cannot be determined.
    \item The rising ratio is a concern the supervisor correctly identifies:
    it is possible the model converges to a suboptimal T2 mAP because the
    KD loss prevents the student from fully adapting.
\end{itemize}

\noindent\textbf{Cost estimate for ablation:} Two additional Stage~2 runs
($\lambda\,=\,0$ and $\lambda\,=\,0.33$), each $\approx$4$\times$H100
for $\approx$18\,h $\approx$~13,824~SBU per run. Total: $\approx$27,648~SBU.
Budget remaining: $\approx$6,859~SBU. \textbf{Not feasible.}

\vspace{0.5em}
\noindent\textbf{Thesis update.} Add to limitations: ``A $\lambda$ ablation
(0, 0.33, 1.0) would be required to confirm whether $\lambda_{\mathrm{kd}}\,=\,1.0$
represents an optimal balance or over-constrains T2 task adaptation. The
increasing kd/det ratio (1.36$\times$$\to$2.44$\times$) suggests distillation
dominates optimisation at later epochs, potentially limiting T2 mAP. This
ablation is left as future work.''

% ----
\subsubsection*{4.5 Tiny AP $\approx$ 0 in Stage~2 --- Forgetting or Unlearned? (Point~5)}

\textbf{Status: Fully answered from existing data --- no new experiments needed.}

\vspace{0.5em}
\noindent T1 (Anti-UAV-RGBT val) tiny-stratum mAP by stage:

\begin{table}[h]
\centering
\caption{Tiny-stratum mAP@0.5 across all stages.}
\small
\begin{tabular}{lrl}
\hline
\textbf{Stage} & \textbf{Tiny mAP@0.5} & \textbf{Note} \\
\hline
After Stage~1 & 0.009 & Never learned well --- only 197/60,620 val frames are tiny \\
After Stage~2 & $\approx$0.001 & Small decline from a near-zero baseline \\
After Stage~3 & 0.0004 & Near zero --- consistent with baseline \\
\hline
\end{tabular}
\end{table}

\noindent\textbf{Interpretation.} The model never achieved meaningful
tiny-target AP in Stage~1 (mAP\,=\,0.009), because Anti-UAV-RGBT contains
only 197 tiny-stratum frames out of 60,620 UAV-present val frames (0.3\,\%).
The supervisor's concern is valid: the Stage~3 tiny ``collapse'' (0.009
$\to$ 0.0004) is not forgetting of a learned capability; it is near-zero
before and after. The FM formula ($\mathrm{FM} = \mathrm{mAP}_{T1} -
0.6725$) uses \textit{overall} mAP, not per-stratum mAP. The overall FM
is dominated by the normal and large strata, where true forgetting occurs.

\vspace{0.5em}
\noindent\textbf{What the data does show for tiny:}
The tiny stratum is a \textit{pre-existing detection failure}, not a
consequence of continual learning. It is correctly excluded from the
gradient starvation argument. The supervisor's point is a valid catch ---
the thesis should be explicit.

\vspace{0.5em}
\noindent\textbf{Thesis update.} Add to the per-stratum discussion in
Chapter~4: ``The tiny stratum (mAP\,=\,0.009 after Stage~1) was never
reliably detected, reflecting the rarity of tiny-stratum instances in
Anti-UAV-RGBT (197 of 60,620 val frames, 0.3\,\%). Stage~3 tiny collapse
(0.009$\to$0.0004) reflects a pre-existing detection failure rather than
catastrophic forgetting of a learned capability. The gradient starvation
argument applies specifically to the normal and large strata, which were
reliably detected after Stage~1 (mAP\,=\,0.719 and 0.461 respectively)
and collapsed substantially after Stage~3.''

% ----
\subsubsection*{4.6 Separating Gradient Starvation from Broader Distribution Shift (Point~6)}

\textbf{Status: Genuine limitation --- goes in limitations.}

\vspace{0.5em}
\noindent CST Anti-UAV differs from Anti-UAV-RGBT in multiple ways simultaneously:
\begin{itemize}
    \item Scale distribution (97.7\,\% tiny/small vs.\ $\approx$15\,\% tiny/small)
    \item Background statistics (indoor/outdoor controlled vs.\ outdoor flight)
    \item Thermal characteristics (possibly different sensor models, acquisition conditions)
    \item Annotation noise (different annotation protocols across datasets)
    \item Camera motion (static camera in CST vs.\ handheld in RGBT)
\end{itemize}

\noindent All of these co-occur in the Stage~3 transition. The gradient starvation
hypothesis attributes forgetting to the scale distribution gap. The supervisor
correctly notes that the evidence is circumstantial.

\vspace{0.5em}
\noindent\textbf{What would distinguish them:}
The clean ablation would be to train Stage~3 on CST filtered to normal+large
targets only (preserving scale distribution) and measure whether T1 normal/large
forgetting is reduced. This would isolate scale-gradient starvation from domain
shift. CST has 0\% large targets, making a large-target-only control impossible;
but a normal-only condition could in principle be constructed.

\vspace{0.5em}
\noindent\textbf{Cost estimate:} Filter CST to normal targets ($\approx$2.3\,\%
of frames), retrain Stage~3 on this subset, evaluate T1 normal forgetting.
Approximately 4$\times$H100, $\approx$10--15\,h $\approx$ 7,680--11,520~SBU.
\textbf{Not feasible.}

\vspace{0.5em}
\noindent\textbf{Best available argument from existing data:}
\begin{enumerate}
    \item The collapse is \textit{stratum-specific} and matches exactly the
    scale gap between datasets. Large-target mAP collapses to 0.000; small-target
    mAP (where CST provides gradient) remains non-zero (0.067). If background/thermal
    distribution shift were the primary driver, forgetting would be expected to be
    more uniform across strata rather than concentrated in the absent stratum.
    \item The weight change is minimal and uniform across layers (per-layer drift
    analysis, Section~4.1). Distribution shift typically drives larger weight
    changes; the high cosine similarity is more consistent with gradient absence
    than with active domain interference.
\end{enumerate}

\noindent\textbf{Thesis update.} Add to limitations: ``The gradient starvation
hypothesis cannot be fully isolated from other confounding factors in the
Stage~2$\to$3 transition, including thermal appearance characteristics,
background statistics, and annotation protocol differences. The control
experiment --- Stage~3 training on CST normal+large targets only --- would
require $\approx$7,680~SBU and is left as future work. The stratum-specificity
of the collapse (large-target mAP\,$\to$\,0.000 while small-target mAP remains
0.067) and the minimal weight change across all layers (mean cosine sim~0.968)
are consistent with gradient starvation as the primary mechanism, but
corroborating evidence from direct gradient measurement or a scale-controlled
ablation is needed for a definitive conclusion.''

% ------------------------------------------------------------
\subsection*{5. Summary Table --- Response Status}

\begin{table}[h]
\centering
\caption{Supervisor feedback response status.}
\small
\begin{tabular}{p{0.04\linewidth}p{0.32\linewidth}p{0.16\linewidth}p{0.38\linewidth}}
\hline
\textbf{\#} & \textbf{Point} & \textbf{Status} & \textbf{Action} \\
\hline
1 & Cosine sim $\neq$ gradient starvation &
    Partially addressed &
    Per-layer breakdown added; structural argument strengthened;
    direct gradient measurement in future work. \\[4pt]
2 & Calibration vs.\ forgetting &
    Fully answered &
    Large AP\,=\,0 rules out calibration.
    Add PR curve figure and one sentence. \\[4pt]
3 & Normal collapses despite CST normal targets &
    Addressed by argument &
    Gradient imbalance + neck feature shift.
    Add to thesis discussion. \\[4pt]
4 & $\lambda$ ablation missing &
    Genuine gap / limitations &
    Not feasible ($\approx$27,648~SBU needed).
    Acknowledge in limitations. \\[4pt]
5 & Tiny AP\,=\,0 in Stage~2 &
    Fully answered &
    Tiny never learned (0.009 after S1).
    Add one paragraph clarifying per-stratum forgetting. \\[4pt]
6 & Cannot separate starvation from domain shift &
    Genuine gap / limitations &
    Not feasible ($\approx$7,680~SBU needed).
    Best available argument documented. Acknowledge in limitations. \\
\hline
\end{tabular}
\end{table}

% ------------------------------------------------------------
\subsection*{6. New Analysis Outputs Available}

\begin{table}[h]
\centering
\caption{New figures and data produced this session.}
\small
\begin{tabular}{p{0.40\linewidth}p{0.52\linewidth}}
\hline
\textbf{File} & \textbf{Use} \\
\hline
\texttt{analysis/stage\_3/drift\_s2s3/fig\_drift\_groups.png}
& Per-layer-group cosine sim S2$\to$S3 --- addresses Point~1. \\
\texttt{analysis/stage\_3/drift\_s2s3/fig\_drift\_cos.png}
& Per-layer cosine similarity dot plot. \\
\texttt{analysis/stage\_3/drift\_s2s3/fig\_drift\_l2.png}
& Per-layer relative L2 drift dot plot. \\
\texttt{analysis/stage\_3/drift\_s2s3/drift\_stats.csv}
& Raw per-layer numbers (335 params). \\
\texttt{analysis/stage\_3/stage3\_naive/eval\_t1/fig\_pr\_curve.png}
& Full PR curve on T1 val after Stage~3 --- addresses Point~2. \\
\texttt{analysis/stage\_3/stage3\_naive/eval\_t1/fig\_conf\_curve.png}
& P/R/F1 vs.\ confidence threshold --- addresses Point~2. \\
\texttt{analysis/stage\_3/stage3\_naive/eval\_t1/scale\_metrics.csv}
& Per-stratum mAP numbers --- addresses Points~2, 3, 5. \\
\hline
\end{tabular}
\end{table}

% ------------------------------------------------------------
\subsection*{7. Thesis Updates Required}

Based on the above analysis, the following additions/edits are needed
in the thesis before the next supervisor response:

\begin{enumerate}
    \item \textbf{Chapter~4 Results --- add per-layer drift table} (Section~4.1
    of this report). Show backbone/neck/head/other cosine sim for S2$\to$S3,
    alongside S1$\to$S2 for comparison. Add \texttt{fig\_drift\_groups.png}.

    \item \textbf{Chapter~4 Results --- add PR curve clarification} (Section~4.2).
    One sentence explaining AP\,=\,0 $\Rightarrow$ recall\,=\,0 at all thresholds
    $\Rightarrow$ not calibration failure. Add \texttt{fig\_conf\_curve.png} in
    appendix.

    \item \textbf{Chapter~4 Results --- clarify tiny stratum} (Section~4.5).
    Explicitly state that tiny-stratum AP was near-zero after Stage~1 (0.009)
    and Stage~3 tiny collapse is not meaningful forgetting.

    \item \textbf{Chapter~4/5 --- add gradient imbalance argument for normal}
    (Section~4.3). 2.3\,\% normal in CST cannot preserve representations tuned
    on 60\,\% normal in Stage~1.

    \item \textbf{Chapter~5 Limitations --- add $\lambda$ ablation gap}
    (Section~4.4). Acknowledge the kd/det ratio observation and the missing
    controlled ablation.

    \item \textbf{Chapter~5 Limitations --- add distribution shift separation gap}
    (Section~4.6). Acknowledge cannot fully isolate gradient starvation from
    thermal/background domain shift. Describe the control experiment that would
    be needed.
\end{enumerate}

% ------------------------------------------------------------
\subsection*{8. Key Numbers for Reference}

\begin{table}[h]
\centering
\caption{All key numbers confirmed as of 16~June~2026.}
\small
\begin{tabular}{llr}
\hline
\textbf{Stage} & \textbf{Metric} & \textbf{Value} \\
\hline
S1 & T1 mAP@0.5 (ceiling) & 0.6725 \\
S1 & Tiny-stratum mAP@0.5 & 0.009 \\
S1 & Normal-stratum mAP@0.5 & 0.719 \\
S1 & Large-stratum mAP@0.5 & 0.461 \\
\hline
S2 & FM (3 seeds) & $-$0.033\,$\pm$\,0.004 \\
S2 & T1 retained & 0.640 \\
S2 & Mean cosine sim S1$\to$S2 & 0.911 \\
S2 & kd/det ratio at best epoch & 2.44$\times$ \\
\hline
S3 naive & FM & $-$0.605 \\
S3 naive & T1 mAP@0.5 (post-training eval) & 0.0761 \\
S3 naive & Large-stratum T1 mAP@0.5 & 0.0000 \\
S3 naive & Normal-stratum T1 mAP@0.5 & 0.0878 \\
S3 naive & Small-stratum T1 mAP@0.5 & 0.0670 \\
S3 naive & Tiny-stratum T1 mAP@0.5 & 0.0004 \\
S3 naive & Mean cosine sim S2$\to$S3 (global) & 0.9676 \\
S3 naive & Cosine sim backbone (learnable) & 0.9951 \\
S3 naive & Cosine sim neck (learnable) & 0.9772 \\
S3 naive & Cosine sim head (learnable) & 0.9894 \\
\hline
\end{tabular}
\end{table}

